%=======================================================
\chapter{Imagens}

\textsc{Qual o pacote?} \textbf{graphicx} (Acesso: \url{https://www.ctan.org/pkg/graphicx})

\textsc{O que faz?}: Insere imagens nos formatos JPG, PNG, PDF ou EPS. Demonstro somente como incluir uma imagem centralizada na página.
\ \\

\begin{codex}{Código para inserção de imagens}
\begin{figure}[h!]
\begin{center}
\includegraphics[scale=.30]{./img/github.png}
\caption{Imagem GitHub}
\label{fig:github}
\end{center}
\end{figure}
\end{codex}
\ \\

\textsc{Resultado:}

\begin{figure}[h!]
	\begin{center}
		\includegraphics[scale=.30]{./img/github.png}
		\caption{Imagem GitHub}
		\label{fig:github}
	\end{center}
\end{figure}

Uma demonstração combinada com o uso do pacote \textbf{multicol} dispondo as figuras lado a lado.
\ \\

\begin{codex}{Código para figuras lado a lado com multicol}
\begin{multicols}{2}
\includegraphics[scale=.15]{./img/github.png}

\includegraphics[scale=.040]{./img/tex-lion.png}
\end{multicols}
\end{codex}
\ \\

\textsc{Resultado:}

\begin{multicols}{2}
	\includegraphics[scale=.15]{./img/github.png}
	
	\includegraphics[scale=.040]{./img/tex-lion.png}
\end{multicols}

Outras opções, tais como imagens lado a lado, aumentar ou diminuir o tamanho, inserir imagem dentro de caixas de texto, etc, são indicadas no pacote.
